\chapter{v6 在备份流水线中的位置}
\label{ch:journey}

\begin{analogybox}[一句话]
v6 不是新的压缩算法，而是\textbf{新的「切香肠刀法」}：同一台 LZ4 压缩机，不同的 chunk 边界；刀法要够快，切出来的块还要能 dedup。
\end{analogybox}

\section{从文件到 chunk 的旅程}

\begin{figure}[htbp]
\centering
\resizebox{\linewidth}{!}{%
\begin{tikzpicture}[node distance=0.45cm]
  \node[streamblock, minimum width=2.0cm] (file) {源文件\\D:\textbackslash CUDA};
  \node[v6block, minimum width=2.2cm, right=0.55cm of file] (cdc) {2F-Gear\\切点};
  \node[block=CoreGreen, minimum width=2.0cm, right=0.55cm of cdc] (hash) {ContentHash};
  \node[block=CorePurple, minimum width=2.0cm, right=0.55cm of hash] (enc) {LZ4/Auto\\压缩};
  \node[block=CoreBlue, minimum width=2.0cm, right=0.55cm of enc] (store) {ChunkStore\\EbPack};
  \draw[arr] (file) -- (cdc);
  \draw[arr] (cdc) -- (hash);
  \draw[arr] (hash) -- (enc);
  \draw[arr] (enc) -- (store);
  \node[lab, below=0.5cm of cdc] {本章焦点};
  \node[lab, below=0.5cm of enc] {与 Stream 相同};
\end{tikzpicture}%
}
\caption{v6 仅替换 CDC 阶段；压缩率对比必须在相同 encode 路径下解读。}
\label{fig:pipeline-journey}
\end{figure}

\section{三代 G-TCDC 的「扫描复杂度」}

\begin{figure}[htbp]
\centering
\begin{tikzpicture}
\begin{axis}[
  ybar,
  width=12cm,
  height=6.5cm,
  bar width=18pt,
  ymin=0, ymax=4.5,
  ylabel={相对 probe 成本（示意）},
  symbolic x coords={Rabin v2, Gear v3, AN v5, {2F v6}, FastCDC Stream},
  xtick=data,
  x tick label style={rotate=25, anchor=east, font=\small},
  legend style={at={(0.5,-0.22)}, anchor=north, legend columns=2, font=\small},
  nodes near coords, nodes near coords style={font=\tiny},
  enlarge x limits=0.12
]
\addplot[fill=CoreRed!60] coordinates {
  (Rabin v2, 4.0)
  (Gear v3, 1.0)
  (AN v5, 3.2)
  ({2F v6}, 1.15)
  (FastCDC Stream, 1.0)
};
\end{axis}
\end{tikzpicture}
\caption{AN v5 因每字节 \texttt{SelectPhaseMask+MixSeedHash} 显著贵于 v6；v6 设计目标即贴近 Stream（实测 per-probe ratio $\approx$1.0--1.2）。}
\label{fig:probe-cost-bar}
\end{figure}

\section{设计约束三角}

\begin{figure}[htbp]
\centering
\begin{tikzpicture}[scale=1.05]
  \node[v6block, minimum width=3.2cm] (perf) at (0,2.8) {性能 $\approx$ Stream};
  \node[v6block, minimum width=3.2cm] (dedup) at (-3.2,-1.2) {独立切点\\opt-in dedup};
  \node[v6block, minimum width=3.2cm] (comp) at (3.2,-1.2) {压缩率\\不劣化};
  \draw[arr] (perf) -- (dedup);
  \draw[arr] (dedup) -- (comp);
  \draw[arr] (comp) -- (perf);
  \node[v6block, minimum width=4cm] at (0,0.3) {v6 2F-Gear\\双域 AND};
\end{tikzpicture}
\caption{v6 在三角中心用 $O(1)$/probe 双指纹达成三目标；v5 偏 dedup 质量牺牲 perf。}
\label{fig:triangle}
\end{figure}
